% Generated draft: requires mathematical and source review.
\documentclass[11pt]{article}
\usepackage[margin=25mm]{geometry}
\usepackage{fontspec}
\setmainfont{Latin Modern Roman}
\usepackage{amsmath,amssymb,mathtools,unicode-math}
\setmathfont{Latin Modern Math}
\usepackage{enumitem,needspace,xurl,hyperref}
\setlength{\parindent}{0pt}
\setlength{\parskip}{5pt}
\setlength{\emergencystretch}{3em}
\newcommand{\N}{\mathbb{N}}
\newcommand{\Z}{\mathbb{Z}}
\newcommand{\Q}{\mathbb{Q}}
\newcommand{\R}{\mathbb{R}}
\newcommand{\C}{\mathbb{C}}
\newcommand{\F}{\mathbb{F}}
\newcommand{\A}{\mathbb{A}}
\newcommand{\PP}{\mathbb{P}}
\newcommand{\E}{\mathbb{E}}
\newcommand{\eps}{\varepsilon}
\newcommand{\dd}{\,\mathrm{d}}
\begin{document}
\section*{UnsolvedMath ID: 20000245}
\textbf{Source problem code:} catalogue entry
\subsection{Definitions and mathematical statement}
Jean Fasel asks: for a smooth affine $k$-variety $X$ of odd dimension $d$ and a rank-$d$ vector bundle $P$, is
\[
c_d(P)=0\text{ in }\operatorname{CH}^d(X)
\quad\Longleftrightarrow\quad
P\simeq\mathcal O_X\oplus Q
\]
for a rank-$(d-1)$ bundle $Q$? Here $c_d(P)$ is the top Chern class. On an affine scheme a free rank-one summand is equivalent to a section with a retraction to $\mathcal O_X$. The splitting implies Chern-class vanishing; the substantive question is the converse, uniformly over the intended fields and odd dimensions. No field restriction appears in the original. The cleaned wording is supplied but remains labeled reconstructed/unverified because of extraction damage.

The attached unreviewed aid reports that the arbitrary-field question remains open, while giving affirmative cases over algebraically closed fields, real closed fields, and characteristic-not-two fields of $2$-cohomological dimension at most one, as well as unspecified additional special cases. Exact theorem hypotheses and citations are absent.

Its proposed reduction assumes $k$ perfect of characteristic not two, $X$ smooth affine of pure odd dimension $d$, and $L=(\det P)^\vee$. Let
\[
e(P)\in\widetilde{\operatorname{CH}}^d(X,L),\qquad
F_L:\widetilde{\operatorname{CH}}^d(X,L)\to\operatorname{CH}^d(X)
\]
be the refined Euler class and forgetful map, with $F_L(e(P))=c_d(P)$. The aid claims
\[
c_d(P)=0\ \Longrightarrow\
e(P)\in\ker(F_L)[2],\qquad\langle-1\rangle e(P)=e(P),
\]
where $A[2]=\{a\in A:2a=0\}$ and $\langle-1\rangle$ is the action of the one-dimensional quadratic form on the twisted Chow--Witt group. For a fixed determinant twist $L$, an Euler-realizable class means an actual $e(P)$ for a rank-$d$ bundle equipped with $\det(P)^\vee\simeq L$. Using an Euler-class splitting criterion with its proper hypotheses, the aid identifies top-Chern detection with absence of nonzero Euler-realizable $2$-torsion in $\ker F_L$. It does not require all ambient kernel torsion to vanish. The precise Euler splitting theorem, dimension range and possible extra field hypotheses must be supplied before that equivalence is used unrestrictedly.
\subsection{Short English statement}
Ask whether vanishing of the top Chern class forces a rank-equals-dimension bundle to split off a trivial line in odd dimension. The aid reduces the issue to realizable refined Euler obstructions, preserving its perfect-field and characteristic restrictions.
\subsection{Sources}
\begin{itemize}
\item Jean Fasel, Question 8 in AIM \emph{Projective modules and $\A^1$-homotopy theory}: \url{https://aimath.org/pastworkshops/projectiveA1problems.pdf}. Canonical location: \texttt{aim-algebraic-geometry-notes.json}, notes[244]. The supplied clean text is reconstructed/unverified; the verification points to page 3 but does not reproduce a full source comparison. No PDF was consulted.
\item The attached aid reports field-dependent known cases and proposes the odd-rank Euler-realizability reduction with low novelty confidence. It is machine generated and not peer reviewed. Original remarks and literature are empty, and no complete authoritative citations are supplied.
\item Local registry: \texttt{data/unsolvedmath/problems.json}, permanent catalogue entry, exact source code \texttt{catalogue entry}.
\end{itemize}
\end{document}
